Here is the updated LaTeX section text with all relevant numbers replaced by their new values:

```latex
\section{Simulation Results}
The simulation results for different governance rules and quorum percentages are summarized below.

For a majority rule with a 4\% quorum, we observe:
\begin{itemize}
    \item The average proposal pass rate is $0.87517$.
    \item Average turnout is $34.92$\%, with a standard deviation of $2.23$\%.
    \item Final Gini coefficient for token distribution is $0.40896$, indicating moderate inequality.
\end{itemize}

For quadratic voting with an 8\% quorum, the results are:
\begin{itemize}
    \item The average proposal pass rate drops to $0.87446$.
    \item Average turnout is slightly higher at $35.07$\%, with a standard deviation of $2.46$\%.
    \item Final Gini coefficient for token distribution increases to $0.40339$, indicating similar levels of inequality as the majority rule but with more variability.
\end{itemize}

For conviction voting with a 4\% quorum, we observe:
\begin{itemize}
    \item The average proposal pass rate is slightly higher at $0.87812$.
    \item Average turnout remains around $34.61$\%, with a standard deviation of $2.15$\%.
    \item Final Gini coefficient for token distribution increases to $0.41388$, indicating somewhat greater inequality in the system compared to majority and quadratic voting rules.
\end{itemize}

In all cases, we observe that higher quorum requirements lead to slightly lower proposal pass rates but also tend to increase average turnout.

The final treasury balance is around $\$11,406.2$ for most scenarios with a standard deviation of approximately \$578. The volatility in the treasury balance ranges from $23.5\%$ to $23.7\%$, indicating relatively stable funding levels across different governance rules and quorum percentages.

In summary, while there are slight differences between the various voting mechanisms and quorum thresholds, all configurations exhibit similar overall performance metrics with respect to proposal pass rates, turnout, inequality measures (Gini coefficients), and treasury stability.
```